\chapter{Conclusion and Future Work}
\label{ch:conclusion}

\section{Conclusion}
\label{sec:conclusion}

This thesis set out to answer whether a clinically meaningful kidney-risk
stratification can be produced from the measurements a routine general checkup
already yields---without the kidney-specific confirmatory tests whose ordering
is the actual bottleneck to early detection---and whether the resulting risk
assessment can be communicated to a lay reader by a large language model whose
output is guaranteed never to diagnose, prescribe, or withhold urgent care.

On the first question, the answer is a qualified yes. A two-stage cascade
trained on 4{,}719 NHANES 2021--2023 participants using fourteen routine inputs
achieves, on a held-out test set of 833 participants, an accuracy of 70.9\%, a
balanced accuracy of 66.0\%, a macro ROC AUC of 0.79, and a high-risk recall of
73.7\%, with only 2 of 38 high-risk participants assigned to the low category.
Set against a majority-class baseline that would achieve higher raw accuracy
while detecting nobody, this represents substantial recoverable signal. The
qualification is that precision on the at-risk categories is low---0.242 for
moderate and 0.394 for high---which is the deliberate and openly reported price
of a gate tuned so that individuals who need a confirmatory test are not
filtered out before they can be referred to one.

On the second question, the answer is more definite, and the mechanism
responsible is not the one that might have been expected. The six-arm ablation
established that between 6.3\% and 10.3\% of raw language-model outputs contain
a critical safety violation in every configuration tested, including the
fine-tuned ones. Neither explicit safety prompting, nor retrieval grounding, nor
fine-tuning on 540 safety-screened examples eliminated this. The deterministic
post-generation validator did: it reduced the delivered critical-unsafe rate to
zero and raised emergency-redirect recall from 65--75\% to 100\%, and these were
the only effects in the entire evaluation to survive multiplicity correction,
doing so at $p<0.001$ for both the base and the fine-tuned model. The safety of
the delivered product is a property of the pipeline, not of the model.

Fine-tuning nonetheless earned its place, on different grounds than safety. It
took format adherence under the compact prompt from 0\% to 100\%, which is what
makes a 25\% reduction in prompt tokens possible; it reduced destructive
fallbacks ninefold, so that patients receive a personalised answer rather than a
generic template far more often; it improved readability by roughly half a
grade level; and, when served with the prompt it was trained against, it
achieved 99.7\% grounding. An unanticipated finding---that serving the adapter
with a different prompt collapses grounding to 67--69\%---yields a practical rule
for deploying fine-tuned models that the implementation now enforces
automatically.

Beyond the two components, the work contributes a factor-aware retrieval
strategy that gives two individuals in the same risk category different
guideline grounding drawn from their own dominant risk factors; a deterministic
citation layer that makes fabricated references structurally impossible; and a
complete, auditable implementation in which every intermediate artefact---stage
probabilities, attributed factors, retrieved passages, raw pre-validator text
and validator decisions---is exposed for inspection.

Finally, the evaluation methodology is itself a contribution. Reporting the
deployed system alone would have produced a headline of ``0\% unsafe output'' and
invited the inference that the fine-tuned model is safe. Only the
component-isolating ablation revealed that the model is not safe and that a
78-line rule system is carrying the guarantee. For safety-critical applications
of generative models, evaluations that cannot attribute the observed behaviour
to a specific mechanism are of limited value, however favourable their headline
numbers.

\section{Future Work}
\label{sec:future}

Eight directions follow from the limitations identified in
Chapter~\ref{ch:discussion}, ordered by the strength of the case for pursuing
them.

\textbf{1. Clinician review of delivered output.} The single most valuable next
step. Both the evaluation detectors and the deployed validator are rule systems
written by the same authors, so a violation neither anticipates is neither
blocked nor counted. Independent adjudication of a stratified sample of
delivered outputs by practising nephrologists or general physicians would place
an empirical bound on the false-negative rate of the guardrail, which no
internal measurement can do.

\textbf{2. Local validation and recalibration.} The model is trained on a United
States survey population and is untested elsewhere. Validation on a Bangladeshi
cohort, with recalibration of the two thresholds against local prevalence and
local costs of referral, is a prerequisite for any deployment in this
institution's context.

\textbf{3. A controlled cascade-versus-single-model comparison.} The
architectural argument for the cascade is sound, but this work did not train a
single three-class classifier on the identical features and splits to test
whether the cascade actually reduces the high-to-low error. That experiment
would convert a design rationale into an empirical claim.

\textbf{4. Closing the emergency-symptom pathway.} The validator's symptom
vocabulary is implemented and tested but unreachable in production, because the
input schema exposes no free-text symptom field. Adding that field end to
end---frontend, schema, contract, validator context---would make an already-built
safety capability live.

\textbf{5. Improving the moderate class.} At 0.505 recall, the moderate category
is the weakest part of the model and the greatest source of headroom. Options
include ordinal-aware loss functions that penalise distant errors more than
adjacent ones, cost-sensitive learning with an explicit clinical cost matrix,
and calibrated probability outputs that would allow the moderate band to be
expressed as an interval rather than a hard category.

\textbf{6. A hybrid safety layer.} The rule-based validator provides a bounded
guarantee for known patterns and cannot generalise. A language-model judge
generalises to novel phrasings but reintroduces probabilistic behaviour. A
layered design---deterministic rules as the non-negotiable floor, with a judge
above it flagging suspicious output that no rule matched, for offline review and
rule extension---would combine the guarantee with a discovery mechanism for the
patterns the rules do not yet cover.

\textbf{7. Longitudinal extension.} Both the data and the system are
cross-sectional. Serial screenings would allow risk trajectory rather than a
point estimate, which is closer to how kidney disease is actually managed and
would make the follow-up guidance substantially more actionable.

\textbf{8. Multilingual education.} The knowledge base and prompts are English
only. For the intended deployment population, Bengali-language education is not
an enhancement but a requirement. BGE-M3 is already multilingual, so the
retrieval layer would transfer; the generation and validation layers would both
need adaptation, and the validator's rule patterns would need to be rewritten
per language---a non-trivial undertaking, since the safety guarantee is only as
good as the language coverage of its rules.

\section{Closing Remarks}
\label{sec:closing}

The system described in this thesis does not diagnose chronic kidney disease and
does not replace a clinician, a creatinine assay, or a urine albumin test. What
it does is convert measurements that a routine checkup already produces into a
risk category with a stated reliability, and then convert that category into
personalised, guideline-grounded, readable education that has passed a
deterministic safety gate before it reaches the reader.

The most transferable lesson concerns where safety comes from. It is tempting,
given a fluent and well-behaved model, to conclude that the model is safe. The
measurements reported here say otherwise: the same model that produced zero
unsafe delivered outputs produced unsafe raw output in one case out of every ten
to sixteen. The difference between those two numbers is not the model. It is the
guardrail---and the only reason that is visible at all is that the evaluation was
designed to look.
